%*******************************************************************************
%*********************************** Chapter 5 *********************************
%*******************************************************************************

\chapter{Conclusions and Recommendations}

This chapter synthesizes the outcomes of the study, presents final conclusions drawn from experimental findings and performance evaluations, and offers practical recommendations for future engineering enhancements and farm-scale deployments.

\section{Summary of Findings}

The design, development, and validation of the integrated automated feeding and watering system with centralized data monitoring yielded the following key findings:

\begin{enumerate}[label=\arabic*.]
    \item \textbf{Dispensing Precision and Automated Scheduling:} The centralized $0.18\text{ m}^3$ hopper and NEMA 17 auger dispensing mechanism achieved a mean Feed Dispensing Accuracy (FDA) exceeding 95\% across 30 repeated trials. Automated feeding schedules were dispatched within $\pm 4.2$ seconds of scheduled target times without manual intervention.
    
    \item \textbf{Edge Sensor Accuracy and Autonomy:} Feeder node instrumentation demonstrated high fidelity. Load cells achieved linear calibration with $R^2 \geq 0.998$ and an error within $\pm 1.2\%$; ultrasonic level sensing maintained depth accuracy within $\pm 0.6$~cm, executing closed-loop water refilling at the 40\% threshold; and the MQ-135 sensor reliably categorized ammonia toxicity into discrete safety tiers.
    
    \item \textbf{Central Processing and Interface Performance:} The Raspberry Pi 4 Model B master unit polled edge nodes over an RS-485/MODBUS RTU industrial serial bus with mean telemetry display latency of $1.14$ seconds. During 72 hours of uninterrupted stress testing, the system achieved 100\% operational uptime with zero database losses and a User Acceptance Score (UAS) of 4.62 / 5.0.
    
    \item \textbf{Multi-Channel Telemetry and Alert Pipeline:} Anomaly detection based on a 14-session rolling window achieved 95\% Alert Response Accuracy (ARA). Remote alert dispatches through the SX1278 LoRa and 4G LTE GSM gateway achieved a 95\% SMS delivery success rate with mean alert latency of $18.4$ seconds.
    
    \item \textbf{Operational Impact:} Time-and-motion comparison indicated a 42.5\% reduction in daily routine caretaker labor requirements, alongside 99.6\% Event-Record Completeness (ERC).
\end{enumerate}

\section{Conclusions}

Based on the empirical findings, the following conclusions are established:

\begin{enumerate}[label=\arabic*.]
    \item The hardwired RS-485/MODBUS RTU multi-node architecture provides an exceptionally stable, noise-immune communication backbone that overcomes the vulnerabilities of conventional wireless implementations in enclosed, metal-framed piggery facilities.
    
    \item Gravimetric closed-loop feed metering combined with ultrasonic liquid level detection successfully eliminates the inaccuracies inherent to volumetric or timed feeding systems, significantly reducing feed waste and preventing livestock dehydration.
    
    \item Integrating environmental gas sensing with behavioral consumption analysis establishes an effective early-warning pipeline for swine illness, empowering caretakers with actionable, timestamped data before visual clinical symptoms manifest.
    
    \item The local touchscreen interface and dual-channel alert system provide an intuitive, high-confidence monitoring ecosystem that bridges advanced embedded systems engineering with practical livestock farm operations in the Philippine agricultural sector.
\end{enumerate}

\section{Recommendations}

To extend the utility, scalability, and technical depth of this system, the following recommendations are suggested for future engineering iterations:

\begin{enumerate}[label=\arabic*.]
    \item \textbf{Individual Animal Identification (RFID Integration):} Integrate ultra-high-frequency (UHF) ear-tag RFID readers at trough stations to transition from pen-level consumption baselines to individual-animal tracking, especially beneficial in gestating sow and farrowing pens.
    
    \item \textbf{Edge Machine Learning Models:} Deploy lightweight embedded machine learning models (such as TinyML on the ESP32-S3 or Raspberry Pi) to combine audio cough detection with consumption and temperature metrics, enabling multi-modal early pathogen screening.
    
    \item \textbf{Solar and Battery Backup Architecture:} Incorporate an uninterruptible power supply (UPS) combining solar harvesting with lithium-iron-phosphate (LiFePO4) battery storage to ensure continuous monitoring and dispensing during rural grid power failures.
    
    \item \textbf{Biosecurity Fogging and Automated Misting:} Expand the actuator controller firmware to drive misting nozzles for automatic evaporative pen cooling and disinfectant spraying when ambient temperatures or ammonia levels exceed warning thresholds.
    
    \item \textbf{Extended Multi-Farm Cloud Telemetry:} Implement optional, secure MQTT-over-TLS telemetry synchronization from the local master unit to a centralized cloud analytics portal, allowing multi-site commercial producers and agricultural cooperatives to benchmark performance across regions.
\end{enumerate}
